% !TEX root = ../manuscript.tex

\section{End-to-End Uncertainty Quantification Workflow}

\subsection{Workflow Overview}

The computational pipeline has three stages: stochastic fault-property
prediction with PREDICT, hierarchical sampling of three-dimensional fault
properties, and field-scale multiphase reservoir simulation. Within the
sampling stage, uncertainty is organized into three levels: (1) geologic
scenario uncertainty, (2) within-window PREDICT stochasticity, and (3)
cross-window coupling uncertainty. Distinguishing the pipeline stages from the
uncertainty levels clarifies which quantities are predicted, which are sampled,
and which are propagated through the reservoir model.

\todo{Add an end-to-end workflow figure that tracks inputs, empirical
libraries, sampled $6\times87$ property fields, upscaled multiphase
properties, reservoir simulations, and quantities of interest.}

\subsection{Stochastic Fault-Property Prediction With PREDICT}

For each geology and throw window, PREDICT generates stochastic fine-scale
fault-core architectures conditioned on the local stratigraphy, fault throw,
and geologic parameters. Each accepted realization yields an upscaled joint
permeability vector $(k_{xx},k_{yy},k_{zz})$. We generate 2,000 accepted
realizations for every geology-window combination. A separate convergence
analysis established that 2,000 realizations bring the component-distribution
distances close to the empirical floor between independent 20,000-realization
reference ensembles.

Adjacent layers of the same lithology are collapsed before fault-core
generation. Clay-smear placement uses the cell-union probability-of-smear
rule: each clay source produces an active cell mask, a cell remains clay if
any source occupies it, and overlapping sources are assigned a parent source
after clay presence is established. This ordering prevents competition among
overlapping sources from converting a potentially clay-occupied cell back to
sand.

\todo{Add the governing PREDICT relations, the material-placement algorithm,
the effective-property upscaling equations, and the final validation evidence
for the cell-union probability-of-smear option.}

\subsection{Hierarchical Sampling of Fault Properties}

\subsubsection{Within-window joint permeability states}

Each PREDICT realization is represented by
\begin{equation}
\mathbf{y}^{(m)} = \left(\log_{10} k_{xx}^{(m)},
\log_{10} k_{yy}^{(m)},\log_{10} k_{zz}^{(m)}\right).
\end{equation}
Joint permeability clusters are detected by $k$-medoids in physical
log-permeability space using Euclidean distance with one log unit as the
component scale. Candidate cluster counts are $K=1,\ldots,5$. Multi-cluster
solutions must satisfy a minimum cluster fraction of 5\%, and the valid
solution with the largest mean silhouette is selected. If its silhouette is
below 0.25, the window is treated as single-cluster.

For interpretation, local percentile ranks $r_{xx}$, $r_{yy}$, and $r_{zz}$
are computed within each window. The joint rank score is
\begin{equation}
s^{(m)} = \frac{r_{xx}^{(m)}+r_{yy}^{(m)}+r_{zz}^{(m)}}{3}.
\end{equation}
Clusters are ordered from low to high by their median joint rank score. Low
and high state libraries each contain 20\% of the full window library.
Complete extreme clusters are included first; if the target size would be
exceeded, only the required lowest- or highest-score samples from the boundary
cluster are retained. This cluster-aware quantile construction preserves
multimodal structure while maintaining a fixed state mass.

Each state medoid is an actual PREDICT realization. The local perturbation
pool contains samples nearest to the medoid within the intersection of the
state library and the medoid-containing cluster. Its size is the smaller of
the available candidate count and the larger of 25 samples or 10\% of that
candidate set. The state-wide pool is the complete low or high state library.
Independent samples are drawn from the complete window library. In every
case, $k_{xx}$, $k_{yy}$, and $k_{zz}$ are selected together from one actual
realization.

\subsubsection{Cross-window similarity and coupling cases}

For each geology, pairwise differences between the six joint permeability
distributions are quantified with normalized multivariate energy distance.
The normalization compares between-window distributional separation with the
characteristic within-window spread. Windows form a similar pair when the
normalized distance is no greater than 0.25. Optional bootstrap resampling
tests the stability of this decision but is not required for the production
workflow because the empirical window libraries were independently
convergence tested.

All 203 partitions of six windows are enumerated. A partition is admissible
only when every pair within every non-singleton group satisfies the similarity
threshold, preventing chain-link grouping. Among admissible partitions, the
workflow first selects the smallest group count and then breaks ties using the
smallest aggregate within-group distance. Singleton windows remain independent
because the distributions do not support coupling them to another window.

Ten multiple-window permeability cases are constructed for each geology: two
independent cases, four fault-wide low/high cases using local or state-wide
pools, and four grouped low/high cases. Non-singleton groups share low/high
state assignments, whereas singleton windows remain independently sampled.
When more than two non-singleton groups are present, candidate binary group
splits are scored by the difference between mean opposite-side and same-side
group distance; the highest-scoring split and its complement are retained.

\subsection{Along-Strike Sampling, Exact Replay, and Multiphase Upscaling}

For each of the 10 cases, the prescribed six-window design is sampled
independently at all 87 along-strike locations. Deterministic seeds and the
accepted PREDICT realization identifiers make every selected permeability
vector exactly replayable. Replay reconstructs the corresponding fine-scale
fault-core maps needed to upscale porosity, capillary pressure, and relative
permeability without storing every fine-scale realization during the original
PREDICT campaign.

The rigorous capillary-pressure workflow uses invasion percolation on each
replayed fault-core realization and retains spatially variable saturation
endpoints. The reservoir-input workflow can either preserve every slice-level
capillary-pressure curve or reduce curves within physically distinct entry-
pressure branches using branch-specific medoids. Dynamic relative-permeability
upscaling is performed for a representative saturation-endpoint realization
in each window and mapped consistently to the retained saturation regions.

\todo{Complete the equations and numerical details for porosity,
invasion-percolation capillary pressure, effective irreducible water
saturation, dynamic relative permeability, branch detection, medoid
selection, endpoint normalization, and quality-control criteria.}
